\section{CPU clients}
\label{sec:cpu}

The CPU client is the bitsliced backend at the host's word width. It
matters for three reasons: it is the verification path for every GPU
report, it is the ECC2K-95 end-to-end control, and the campaign's
spot fleet included CPU workers whose points are indistinguishable
from the GPUs'.

\begin{table}[t]
\centering
\small
\caption{Bitsliced CPU client, $\mathrm{GF}(2^{131})$, one core, median
  of five. Speedups are against the g++ 64-lane build on the same host. The 2009 qhasm implementation of Bailey et al.\ reached
  $533$ cycles per iteration on a Core~2 with 128-bit vectors.}
\label{tab:cpu}
\begin{tabular}{@{}llrrl@{}}
\toprule
host & word & lanes & M it/s & note \\
\midrule
Xeon 2.1~GHz & \texttt{uint64\_t} & 64 & 4.42 & 475 cycles/iteration \\
Xeon 2.1~GHz & AVX2 & 256 & 10.7 & 196 cycles/iteration \\
Xeon 2.1~GHz & AVX-512 & 512 & 12.6 & 167 cycles/it.; 4 threads 40~M \\
M4 Pro, g++ 16.1 & NEON & 128 & 11.342 & $2.19\times$ over 64 lanes (5.185) \\
M4 Pro, Apple clang 17 & NEON & 128 & 15.673 & $3.02\times$ over g++/64 \\
Graviton3 (V1, SHA3), clang & NEON & 128 & 6.214 & $2.38\times$ over g++/64 \\
Graviton2 (N1, no SHA3), clang & NEON & 128 & 4.220 & $2.09\times$ over g++/64 \\
Modal container, 4 threads & AVX-512 & 512 & 12.4 & 49.4~M it/s on 4 threads \\
\bottomrule
\end{tabular}
\end{table}

Table~\ref{tab:cpu} is the ledger.\art{ecc2k130/README.md, ``Throughput''}
Two of its facts generalise. The AVX-512 build spends half its
multiply on moves rather than arithmetic ($4{,}069$ logic instructions
against $4{,}287$ \texttt{vmovdqa32}) because the Karatsuba leaf keeps
254 values live against 32 registers; a scheduling pass over the same
DAG recovered moves only in the conversion ($4{,}287$ to $3{,}903$),
missed its declared target, and recorded the leaf-size and compiler
levers as measured dead ends.\art{ecc2k130/HOST-SCHEDULE.md} On
aarch64 the gain splits cleanly: the NEON word with \texttt{BSL} alone
is $1.76\times$ on the M4~Pro, \texttt{EOR3} a further $1.24\times$, and a core without \texttt{FEAT\_SHA3} (Graviton2) still takes
$1.81\times$ from the wider word and \texttt{BSL}. Compiler choice is worth more than either on
this code: gcc emits $9{,}372$ instructions for the multiply leaf,
$59\%$ of them loads and stores against spilled temporaries, where
clang emits $1{,}902$ plus five outlined calls. The SHA3 probe is gated
on the CPU reporting the feature, because an ungated probe on a
Graviton2 selects a \texttt{+sha3} spelling under both compilers and
the resulting binary dies with \texttt{SIGILL} on its $8{,}907$
\texttt{eor3} instructions.\art{ecc2k130/README.md, ``Word width''}

In the fleet a CPU worker at $2.3$~B/s (the self-reported median of
the twelve CPU slots on 2026-09-17) is about a sixth of an
RTX~PRO~6000; on Modal a 32-thread (16-core) sidecar beside a GPU adds about
$25\%$ to the container's cost for $2.7\%$ more iterations and is off
by default.\art{ecc2k130/benchmarks/dp-interval/raw/worker-counters-2026-09-17.tsv}\art{ecc2k130/README.md, ``Walking on the container's CPUs too''}
